% !TeX root = ../navier-stokes-manuscript.tex
% ============================================================
% 2.1 Vortex Stretching
% ============================================================

\subsection{Vortex Stretching}\label{sub:ThePerfectlyStretchingVortex}

The momentum equation has four pieces:
\[
  \underbrace{\partial_t u}_{\text{local acceleration}}
  +
  \underbrace{(u\cdot\nabla)u}_{\text{nonlinear transport}}
  =
  \underbrace{-\nabla p}_{\text{pressure}}
  +
  \underbrace{\nu\Delta u}_{\text{viscous diffusion}}.
\]
The central regularity problem lies in the competition between nonlinear transport and viscosity.
The nonlinear term transports the velocity field by itself, so its own motion can sharpen the gradients it carries.
Viscosity acts on those gradients in the opposite direction and smooths sharp variations.

The hope is simple.
Viscosity smooths the fluid faster than nonlinear transport can sharpen it.
If that remains true at every scale, the velocity stays smooth.
A finite-time singularity would require the sharpening to keep moving into finer structure until some derivative becomes unbounded before viscosity smooths it out.

\divider{Vortex Stretching}

Vortex stretching is the three-dimensional motion that can keep this sharpening alive.
Its physical danger, the thinning it produces, and its vorticity equation are three views of that one mechanism.
Vorticity and the symmetric strain are
\[
  \omega:=\nabla\times u,
  \qquad
  S:=\frac12\bigl(\nabla u+(\nabla u)^{\mathsf T}\bigr).
\]
The vector \(\omega\) measures local rotation, while \(S\) measures extension and contraction in the surrounding velocity field.
In two dimensions, the stretching term vanishes.
In three dimensions, the surrounding velocity can extend a rotating structure along its axis.
If \(\omega\) aligns with a positive extensional direction of \(S\), then
\[
  \omega\cdot S\omega>0,
\]
so the strain amplifies the vorticity.

Incompressibility ties that extension to a contraction:
\[
  \operatorname{tr}S=\nabla\cdot u=0.
\]
Extension in one direction must be accompanied by a net contraction across the transverse cross-section.
Suppose the rotating structure has transverse width \(\ell\) and the velocity changes by \(\Delta U\) across that width.
The corresponding gradient is roughly
\[
  |\nabla u|\sim\frac{\Delta U}{\ell}.
\]
If \(\ell\) shrinks while \(\Delta U\) does not shrink equally fast, the gradient grows.
Thinning by itself guarantees nothing; the velocity change across the shrinking width has to persist.

Finite kinetic energy still allows this concentration.
Energy measures
\[
  \int_{\R^3}|u(x,t)|^2\,dx,
\]
whereas a gradient measures how quickly \(u\) changes across a distance.
A sharp structure can occupy such a small volume that its energy remains finite, or even tends to zero, while its gradients diverge.

The curl of the momentum equation makes the competition exact.
Pressure disappears because
\[
  \nabla\times\nabla p=0,
\]
and the vorticity satisfies
\[
  \partial_t\omega+(u\cdot\nabla)\omega
  =
  (\omega\cdot\nabla)u+\nu\Delta\omega
  =
  S\omega+\nu\Delta\omega.
\]
Here \(S\omega\) is vortex stretching and \(\nu\Delta\omega\) is viscous diffusion.
Pairing the equation with \(\omega\) gives
\[
  \frac12(\partial_t+u\cdot\nabla)|\omega|^2
  =
  \omega\cdot S\omega
  +
  \nu\,\omega\cdot\Delta\omega.
\]
The stretching term records amplification by aligned extension.
The viscous term has no fixed sign at an individual point, since
\[
  \nu\,\omega\cdot\Delta\omega
  =
  \frac{\nu}{2}\Delta|\omega|^2
  -
  \nu|\nabla\omega|^2.
\]
After integration over \(\R^3\), incompressibility removes the transport integral and decay removes the Laplacian flux:
\[
  \frac12\frac{d}{dt}\norm{\omega(t)}_2^2
  +
  \nu\norm{\nabla\omega(t)}_2^2
  =
  \int_{\R^3}\omega\cdot S\omega\,dx.
\]
The dangerous case is persistent aligned extension across shrinking scales, with strain continuing to amplify vorticity while viscosity diffuses the same sharpening structure.

\divider{\NavierStokes{} Scaling}

A shrinking width changes space, velocity, pressure, and time together.
The equation itself fixes that relation.
For \(\lambda>0\), define
\[
  u_\lambda(x,t)
  :=
  \lambda u(\lambda x,\lambda^2t),
  \qquad
  p_\lambda(x,t)
  :=
  \lambda^2p(\lambda x,\lambda^2t).
\]
Each of the four momentum terms receives the same factor:
\[
\begin{aligned}
  \partial_t u_\lambda(x,t)
  &=
  \lambda^3(\partial_t u)(\lambda x,\lambda^2t),
  &
  (u_\lambda\cdot\nabla)u_\lambda(x,t)
  &=
  \lambda^3\bigl((u\cdot\nabla)u\bigr)(\lambda x,\lambda^2t),
  \\
  \nabla p_\lambda(x,t)
  &=
  \lambda^3(\nabla p)(\lambda x,\lambda^2t),
  &
  \Delta u_\lambda(x,t)
  &=
  \lambda^3(\Delta u)(\lambda x,\lambda^2t).
\end{aligned}
\]
Local acceleration, nonlinear transport, pressure, and viscosity all scale together.
The corresponding scales are
\[
  x\mapsto\lambda^{-1}x,
  \qquad
  u\mapsto\lambda u,
  \qquad
  p\mapsto\lambda^2p,
  \qquad
  t\mapsto\lambda^{-2}t.
\]
This scaling compares solutions at different scales.
A cascade inside one fixed solution is an additional dynamical claim.

\divider{Critical Height}

We now need a measure that the \NavierStokes{} scaling neither enlarges nor diminishes.
The homogeneous Sobolev scale identifies it.
Let \(\Lambda:=(-\Delta)^{1/2}\).
For every \(s\in\R\) for which the homogeneous norm is finite,
\[
  \Lambda^s u_\lambda(x,t)
  =
  \lambda^{1+s}(\Lambda^s u)(\lambda x,\lambda^2t).
\]
Squaring the \(L^2\) norm and changing variables from \(x\) to \(y=\lambda x\) gives
\[
\begin{aligned}
  \norm{u_\lambda(t)}_{\dot H^s}^2
  &=
  \int_{\R^3}
  \lambda^{2+2s}
  \bigl|(\Lambda^s u)(\lambda x,\lambda^2t)\bigr|^2
  \,dx
  \\
  &=
  \lambda^{2s-1}
  \norm{u(\lambda^2t)}_{\dot H^s}^2.
\end{aligned}
\]
The exponent vanishes exactly when
\[
  2s-1=0,
  \qquad
  s=\frac12.
\]
This gives the critical height
\[
  H(t)
  :=
  \frac12\norm{u(t)}_{\dot H^{1/2}}^2
  =
  \frac12\norm{\Lambda^{1/2}u(t)}_2^2.
\]

The kinetic-energy identity proved in \Cref{prop:ClassicalKineticEnergyIdentity} supplies the level below the critical height:
\begin{equation}\label{eq:BodyEnergyIdentity}
  \frac12\norm{u(t)}_2^2
  +
  \nu\int_0^t\norm{\nabla u(s)}_2^2\,ds
  =
  \frac12\norm{u_0}_2^2.
\end{equation}
Writing \(H_\lambda\) for the critical height of \(u_\lambda\), the energy, critical, and first-derivative levels scale as
\[
  \norm{u_\lambda(t)}_2^2
  =
  \lambda^{-1}\norm{u(\lambda^2t)}_2^2,
  \qquad
  H_\lambda(t)=H(\lambda^2t),
  \qquad
  \norm{\nabla u_\lambda(t)}_2^2
  =
  \lambda\norm{\nabla u(\lambda^2t)}_2^2.
\]
Under this scaling, the energy level falls, critical height is unchanged, and first-derivative size rises.
The critical height sits between the finite energy and the derivative growth that a concentration scenario would require.

\divider{Nonlinear Production versus Viscous Dissipation}

Scale invariance tells us what \(H\) measures.
\NavierStokes{} tells us how \(H\) changes.
Apply \(\Lambda^{1/2}\) to the momentum equation and pair with \(\Lambda^{1/2}u\).
Define the pressure-complete signed production of critical height by
\[
  \mathcal P_{1/2}(t)
  :=
  -\operatorname{Re}
  \pair
    {\Lambda^{1/2}u(t)}
    {\Lambda^{1/2}\bigl((u(t)\cdot\nabla)u(t)+\nabla p(t)\bigr)}_{L^2}.
\]
The pressure pairing vanishes because \(\nabla\cdot u=0\), while viscosity contributes \(-\nu\norm{\Lambda^{3/2}u}_2^2\).
The critical-height balance is
\[
  H'(t)
  +
  \nu\norm{\Lambda^{3/2}u(t)}_2^2
  =
  \mathcal P_{1/2}(t).
\]
This is the \(n=1\) case of the pressure-complete calculation in \Cref{prop:CriticalHeightSpatialAscent}.
Positive \(\mathcal P_{1/2}\) contributes toward growth, but \(H\) rises only when that production exceeds viscous dissipation.
The production has no assumed sign.
The opening competition has now reached the scale that concentration itself leaves unchanged.

\divider{The Heat Clock}

The viscous term also assigns a time to each active width.
This is exact for the linear heat equation.
If \(v_t=\nu\Delta v\), then
\[
  \widehat v(\xi,t+\delta t)
  =
  e^{-\nu|\xi|^2\delta t}\widehat v(\xi,t).
\]
A scale-localized structure with characteristic width \(r\) has characteristic frequency \(|\xi|\simeq r^{-1}\).
Order-one damping occurs when \(\nu|\xi|^2\delta t\) is of order one, so the viscous comparison time is
\begin{equation}\label{eq:BodyHeatComparisonTime}
  \tau_\nu(r)
  \asymp
  \frac{r^2}{\nu}.
\end{equation}
Halving the width quarters the heat clock.
For a proposed terminal time \(T_*\), the width whose heat clock matches the remaining time satisfies
\[
  r_\nu(t)
  \asymp
  \sqrt{\nu(T_*-t)}.
\]
The calculation attaches a comparison time to a width already present in the flow.
The nonlinear evolution still has to create and occupy that width.

\divider{The Zeno Cascade}

Now place geometrically shrinking widths beside their heat clocks.
Set
\[
  r_n:=2^{-n}r_0,
  \qquad
  \tau_n:=\frac{r_n^2}{\nu}.
\]
Using \eqref{eq:BodyHeatComparisonTime},
\[
  \tau_n
  =
  \frac{r_0^2}{\nu}4^{-n},
  \qquad
  \sum_{n=0}^{\infty}\tau_n
  =
  \frac{4r_0^2}{3\nu}
  <
  \infty.
\]
Infinitely many shrinking comparison clocks fit arithmetically inside a finite total time.
This is the timing of the Zeno cascade.

An actual cascade requires one velocity history to form widths
\[
  r_0>r_1>r_2>\cdots\longrightarrow0
\]
at times
\[
  t_0<t_1<t_2<\cdots\uparrow T_*<\infty,
\]
with \(t_{n+1}-t_n\) following a summable schedule comparable to \(\tau_n\).
Every transition has to occur in the same fluid while viscosity is already acting on the newly formed scale.
The geometric series establishes the arithmetic possibility.
The coupled \NavierStokes{} evolution carries the burden of realizing it.

The general critical-height route assumes no repeating core and no fixed profile.
If the maximal time is finite, the endpoint \(L^3\) criterion of Escauriaza, Seregin, and \v{S}ver\'ak gives \parencite{EscauriazaSereginSverak2003}
\[
  T_*<\infty
  \quad\Longrightarrow\quad
  \limsup_{t\uparrow T_*}\norm{u(t)}_{L^3}=\infty.
\]
Since \(u=\Lambda^{-1/2}\Lambda^{1/2}u\), the Hardy--Littlewood--Sobolev inequality gives \parencite[Theorem~4.3]{LiebLoss2001}
\[
  \norm{u(t)}_{L^3}
  \leq
  C\norm{u(t)}_{\dot H^{1/2}}
  =
  C\sqrt{2H(t)}.
\]
Combined with the terminal \(L^3\) divergence, this produces a sequence \(t_n\uparrow T_*\) for which
\[
  H(t_n)\longrightarrow\infty.
\]
Plancherel and Cauchy--Schwarz reduce the interpolation step to
\[
\begin{aligned}
  2H(t)
  &=
  \norm{u(t)}_{\dot H^{1/2}}^2
  =
  \int_{\R^3}|\xi|\,|\widehat u(\xi,t)|^2\,d\xi
  \\
  &\leq
  \left(\int_{\R^3}|\widehat u(\xi,t)|^2\,d\xi\right)^{1/2}
  \left(\int_{\R^3}|\xi|^2|\widehat u(\xi,t)|^2\,d\xi\right)^{1/2}
  \\
  &=
  \norm{u(t)}_2\norm{\nabla u(t)}_2.
\end{aligned}
\]
Along the terminal sequence, this turns critical-height growth into global derivative growth.
When \(\norm{u_0}_2=0\), \eqref{eq:BodyEnergyIdentity} reads \(\norm{u(t)}_2^2+2\nu\int_0^t\norm{\nabla u(s)}_2^2\,ds=0\), so \(u=0\) throughout the interval and the zero pair continues it for all time; hence the finite-endpoint branch has \(\norm{u_0}_2>0\).
On that branch, the same equation changes the interpolation bound into
\[
  \norm{\nabla u(t_n)}_2
  \geq
  \frac{2H(t_n)}{\norm{u_0}_2}
  \longrightarrow
  \infty.
\]
This general route forces the critical height and the global \(L^2\) size of the first spatial derivatives to become unbounded along a terminal sequence.
It leaves open whether the pointwise velocity amplitude \(\norm{u(t)}_{L^\infty}\) diverges.

The critical spike is a separate fixed-profile model with an explicit amplitude-width law.
Choose a nonzero divergence-free profile \(\phi\in\Ssigma\) and set
\[
  u_\ell(x)
  :=
  \ell^{-1}\phi\!\left(\frac{x}{\ell}\right).
\]
For this ansatz,
\[
\begin{aligned}
  U_\ell
  :=
  \norm{u_\ell}_{L^\infty}
  &\sim
  \ell^{-1}
  \longrightarrow
  \infty,
  &
  \norm{\nabla u_\ell}_{L^\infty}
  &\sim
  \ell^{-2}
  \longrightarrow
  \infty,
  \\
  \tau_\ell
  &\asymp
  \frac{\ell^2}{\nu}
  \longrightarrow
  0,
  &
  E_\ell
  :=
  \frac12\norm{u_\ell}_2^2
  &\sim
  U_\ell^2\ell^3
  \sim
  \ell
  \longrightarrow
  0.
\end{aligned}
\]
The velocity amplitude and gradients diverge while the energy carried by the shrinking profile tends to zero.
At the same time,
\[
  \frac12\norm{u_\ell}_{\dot H^{1/2}}^2
  =
  \frac12\norm{\phi}_{\dot H^{1/2}}^2,
\]
so its critical height stays fixed.
The spike is a parallel amplitude-width model, separate from the route in which \(H(t_n)\) becomes unbounded.

For either picture to describe a finite-time singularity, the same velocity field would have to pass through sharper states on shorter clocks as \(t\uparrow T_*\).
The next question is what that collapsing clock requires of each successive time response of the same field.
